\noindent\textbf{Elección del procedimiento.} La polar es $C_D=C_{D0}+KC_L^2$, con $K=1/(\pi e\mathcal A)$. Para maximizar $C_L/C_D$, derivar respecto a $C_L$ e igualar el numerador a cero: $C_{D0}-KC_L^2=0$. Se elige la raíz positiva.

\begin{enumerate}

\item Obtener el coeficiente cuadrático del alargamiento y Oswald.

\[\fitmath{K=\frac{1}{\pi\cdot e\cdot \mathrm{AR}}}\]

\[\fitmath{K=\frac{1}{\pi\cdot (0{,}82)\cdot (9{,}5)}\approx 0{,}04086135\,}\]

\item Despejar el coeficiente óptimo de la condición de derivada nula.

\[\fitmath{\mathrm{CLopt}=\sqrt{\frac{C_{D0}}{K}}}\]

\[\fitmath{\mathrm{CLopt}=\sqrt{\frac{(0{,}018)}{(0{,}04086134611)}}\approx 0{,}6637124\,}\]

\item Evaluar la eficiencia con el óptimo; equivale a $1/(2\sqrt{C_{D0}K})$.

\[\fitmath{\mathrm{LDmax}=\frac{\mathrm{CLopt}}{C_{D0}+K\cdot {\mathrm{CLopt}}^{2}}}\]

\[\fitmath{\mathrm{LDmax}=\frac{(0{,}6637123789)}{(0{,}018)+(0{,}04086134611)\cdot {(0{,}6637123789)}^{2}}\approx 18{,}43645\,}\]

\end{enumerate}
